\documentclass[12pt]{article}

%%\usepackage[english]{babel}
\usepackage{amsmath}
\usepackage{amssymb}
\usepackage{url}


\begin{document}

\title{Econometrics Notes IV: Binomial Series and Fractional Differencing}
\author{Artur Wegrzyn}
\maketitle

\abstract{This note introduces binomial series and briefly explains how it is used for fractional differencing in time series analysis. }


\tableofcontents

\section{Definition of Binomial Series}
The familiar binomial formula:
\begin{equation*}
(x+y)^n = \sum_{k=0}^{n} \binom{n}{k} x^{n-k} y^{k}
\end{equation*}

with $n \in \mathbb{N}$ is generalized to arbitrary 
real exponent $\alpha$ in the following manner:

\begin{equation}
(1 + x)^{\alpha} = \sum_{k=0}^{\infty} \binom{\alpha}{k} x^k
\end{equation}

This begs the question of how to define the binomial coefficient 
$\binom{\alpha}{k}$ for fractional $\alpha$. \\
\indent Consider first the formula for binomial coefficient:

\begin{eqnarray*}
\binom{n}{k} &=& \frac{n!}{k!(n-k!)} = \frac{n(n-1)\dots (n-k+1)}{k!} = \frac{n^{\underline{k}}}{k!}
\end{eqnarray*}

where falling factorial notation $x^{\underline{n}}$ is used:
\begin{equation*}
x^{\underline{n}} = \prod_{k=1}^{n} x - k + 1
\end{equation*}

Observe that for $n=0$ this yields:
\begin{equation*}
x^{\underline{0}} = \prod_{k=1}^{0}x - k + 1 = 1
\end{equation*}
since value of empty product is one by definition. For $n=0$ and $n=1$ we get:
\begin{eqnarray*}
x^{\underline{1}} &=& x \\
x^{\underline{2}} &=& x (x-1)
\end{eqnarray*}
 

Hence, the generalized coefficient is defined thus:
\begin{equation*}
\binom{\alpha}{k} = 
\frac{\alpha^{\underline{k}}}{k!} = 
\frac{ \alpha (\alpha - 1) \dots (\alpha - k + 1) }{k!}
\end{equation*}

For later reference, the formulae for generalized binomial coefficient for $k = 0, 1, 2, 3$ are:
\begin{eqnarray*}
\binom{\alpha}{0} &=& \frac{\alpha^{\underline{0}}}{0!} = \frac{1}{1} = 1 \\
\binom{\alpha}{1} &=& \frac{\alpha^{\underline{1}}}{1!} = \frac{\alpha}{1} = \alpha \\
\binom{\alpha}{2} &=& \frac{\alpha^{\underline{2}}}{2!} = \frac{\alpha(\alpha - 1)}{2} \\
\binom{\alpha}{3} &=& \frac{\alpha^{\underline{3}}}{3!} = \frac{\alpha(\alpha - 1)(\alpha-2)}{6}
\end{eqnarray*}


\section{Application to Fractional Differencing}
Consider time series $\{ y_t \}_{t \in T}$. The usual first-order differencing is:

\begin{equation*}
\Delta y_t = y_t - y_{t-1} = y_t - L y_t = (1-L) y_t
\end{equation*}

this is generalized to fractional $d$:

\begin{equation*}
\Delta^{d}y_t = (1-L)^d y_t
\end{equation*}

Expanding right-hand side using binomial series formula:
\begin{equation*}
(1-L)^d y_t = \sum_{k=0}^{\infty} \binom{d}{k} (-1)^k L^k
\end{equation*}

Writing out first few elements of the series:
\begin{equation*}
(1-L)^d y_t = 1 - d y_{t-1} - \frac{1}{2}d(1-d)y_{t-2} - \frac{1}{6}(1-d)(2-d)y_{t-3} - \cdots
\end{equation*}


The idea of fractional differencing is discussed more thoroughly in \cite{fractional_differencing_hosking}. Of particular relevance is disussion in chapter 5 of \cite{fuller_intro_tsa} on applications of fractional differencing to analysis of financial time series.

\bibliography{tsa_bibliography}
\bibliographystyle{ieeetr}

\end{document}
